%========================================================
% CHAPTER 1 — General Framework
%========================================================

\chapter{General Framework}

%--------------------------------------------------------
% Section 1.1 — Purpose and Scope
%--------------------------------------------------------

\section{Purpose and Scope}

The purpose of this work is to establish a coherent framework for the mathematical modeling and numerical simulation of polarization-entangled photon pairs propagating through optical fibers.

The central objective is to organize the physical assumptions, mathematical formalism, and simulation strategy into a reproducible structure, so that each computational step can be traced back to a precise theoretical element. In this way, the simulator is not developed as an isolated numerical tool, but as a direct operational translation of the underlying quantum model.

\medskip

The modeling approach adopted throughout this work is intentionally progressive. Simplified analytically tractable models are introduced first, and are then extended toward increasingly realistic descriptions incorporating statistical effects, decoherence, and effective noisy evolution.

This strategy enables a controlled connection between physical interpretation, analytical predictions, and numerical validation.

%--------------------------------------------------------
% Section 1.2 — Physical Scenario and Experimental Setting
%--------------------------------------------------------

\section{Physical Scenario and Experimental Setting}

The physical system considered in this work consists of a source of polarization-entangled photon pairs, two independent optical-fiber arms, and a final polarization-measurement stage.

A spontaneous parametric down-conversion (SPDC) source generates photon pairs ideally prepared in a Bell state. The polarization degree of freedom encodes the qubit associated with each photon, with the horizontal and vertical polarization modes identified as:

\[
\{|H\rangle, |V\rangle\}
\]

After generation, the two photons are separated into two spatially distinct paths, denoted \(A\) and \(B\), and propagate through independent optical fibers.

\medskip

The fibers affect the polarization state of the propagating photons. In the experimental setting considered here, the relevant propagation effects are represented by two main physical mechanisms.

The first one is a residual coherent phase shift between the horizontal and vertical polarization components. After partial polarization compensation, the remaining deterministic or slowly varying effect is modeled as a relative phase accumulated during propagation. Physically, this phase originates from birefringence, optical-path differences, mechanical perturbations, and thermal fluctuations along the fiber. At the two-photon level, this modifies the phase relation between the components of the entangled state and therefore changes the measured polarization correlations.

\medskip

The second mechanism is local depolarization. While the phase shift describes a coherent transformation of the polarization state, depolarization represents an effective loss of polarization information in each fiber arm. This accounts for unresolved coupling to additional degrees of freedom, imperfect compensation, environmental instability, or fluctuations that are not explicitly tracked at the measurement stage. Since the two photons propagate through different fibers, this effect is treated locally and independently on arms \(A\) and \(B\).

\medskip

At the output of the fibers, the photons are collected by polarization analyzers and single-photon detectors. Measurements are performed on the two arms jointly, so that the relevant experimental information is contained in coincidence statistics rather than in single-arm detection probabilities alone.

From these joint measurements one extracts polarization correlations and state-level metrics.

\medskip

The physical scenario can therefore be summarized as:

\[
\begin{array}{c}
\text{SPDC entangled-photon source} \\
\downarrow \\
\text{Separation into fiber arms } A \text{ and } B \\
\downarrow \\
\text{Residual phase shift and local depolarization} \\
\downarrow \\
\text{Polarization analysis and coincidence detection} \\
\downarrow \\
\text{Extraction of correlations and entanglement metrics}.
\end{array}
\]

%--------------------------------------------------------
% Section 1.3 — Theoretical Formalism and Notation
%--------------------------------------------------------

\section{Theoretical Formalism and Notation}

The physical scenario described above is formulated within the language of finite-dimensional quantum information theory. The two photons are treated as a bipartite quantum system, where each subsystem carries one polarization qubit and the joint state belongs to the tensor-product space associated with arms \( A \) and \( B \).

\medskip

The notation follows, as far as possible, the standard conventions of quantum information theory, in particular those used in \textit{Quantum Computation and Quantum Information} by Nielsen and Chuang~\cite{nielsen_chuang_qcqi}. Pure states are denoted by kets such as \( |\psi\rangle \), density operators by \( \rho \), local subsystems by labels \( A \) and \( B \), and tensor products by \( \otimes \).

\medskip

Single-photon polarization states are expressed in the computational basis identified with the horizontal and vertical polarization modes,

\[
|H\rangle \equiv |0\rangle,
\qquad
|V\rangle \equiv |1\rangle .
\]

The corresponding two-photon basis is written as

\[
\{
|00\rangle,
|01\rangle,
|10\rangle,
|11\rangle
\},
\]

where the first label refers to arm \( A \) and the second label refers to arm \( B \).

\medskip

Coherent polarization transformations are described by local operators acting on the individual arms, while effective noisy processes are described at the density-operator level. This allows the same formalism to treat pure-state evolution, statistical mixtures, and non-unitary propagation effects.

\medskip

Observable quantities are expressed through expectation values of Hermitian operators. In particular, polarization correlations are described using Pauli operators acting locally on the two arms.

\medskip

The propagated state is characterized through a fixed set of quantities used throughout the work: fidelity \( F \), purity \( \gamma \), concurrence \( C \), tangle \( \tau \), entanglement of formation \( E_F \), and polarization correlations \( \langle \sigma_i \otimes \sigma_j \rangle \). Their formal definitions are introduced in the dedicated theoretical chapters, while Chapter~1 only establishes the notation used to refer to them.

%--------------------------------------------------------
% Section 1.4 — Computational Architecture and Framework
%--------------------------------------------------------

\section{Computational Architecture and Framework}

The computational framework developed in this work is designed as a modular numerical counterpart of the theoretical description. Its purpose is to make each physical or mathematical object appearing in the text correspond to a clear computational component in the codebase.

\medskip

The simulator is organized into functional layers for state preparation, operator construction, state evolution, quantum channels, measurement evaluation, metric computation, testing, and experiment-specific workflows. This organization separates reusable low-level routines from higher-level numerical experiments, improving reproducibility and making each result traceable to the theoretical object from which it originates.

\medskip

State routines construct Bell states and density matrices; evolution routines apply local unitary transformations; channel routines implement global, local, two-arm, and composite fiber-channel maps; measurement routines evaluate correlation functions, correlation tensors, and CHSH quantities; metric routines compute the scalar indicators used to characterize the propagated state.

\medskip

Numerical experiments are built as higher-level workflows that combine these validated components. Each experiment specifies an input state, a propagation mechanism, a set of observables or metrics, and a comparison with analytical predictions whenever available. This makes the simulator not only a numerical tool, but also a reproducibility layer connecting physical assumptions, mathematical formulation, and numerical results.

\medskip

The role of the computational architecture is therefore to support a one-to-one correspondence between theoretical modeling and numerical implementation. Each output of the simulator should be interpretable in terms of a precise state preparation, propagation mechanism, measurement operator, or metric definition.